% ENEM 2017-2025 — 28 questões MCQ — Análise Combinatória II, Probabilidade I e Probabilidade II
% Fonte: listas "Análise Combinatória II", "Probabilidade I" e "Probabilidade II" do site projetoagathaedu.com.br
% Omitidas em Análise Combinatória II: Q3,Q5,Q6,Q8,Q9,Q10,Q11,Q15,Q19,Q20 (imagens essenciais ou alternativas irrecuperáveis)
% Omitidas em Probabilidade I: Q6,Q8,Q12,Q14,Q15,Q16,Q19,Q20,Q24,Q26,Q27 (imagens essenciais ou alternativas irrecuperáveis)
% Omitidas em Probabilidade II: Q2,Q3,Q5,Q6,Q8,Q10,Q11,Q12,Q14,Q15,Q16,Q19 (imagens essenciais ou alternativas fórmulas corrompidas)

---
tipo: Múltipla Escolha
dificuldade: Fácil
nivel: médio
disciplina: Matemática
assunto: Análise Combinatória
gabarito: D
tags: [aparelho, transmissão, frequência, chaves]
fonte: concurso
concurso: ENEM-PPL
banca: INEP
ano: 2019
numero: 1
---
\question
Uma pessoa comprou um aparelho sem fio para transmitir músicas a partir do seu computador para o rádio de seu quarto. Esse aparelho possui quatro chaves seletoras e cada uma pode estar na posição 0 ou 1. Cada escolha das posições dessas chaves corresponde a uma frequência diferente de transmissão.

A quantidade de frequências diferentes que esse aparelho pode transmitir é determinada por
\begin{oneparchoices}
\choice 6
\choice 8
\choice 12
\correctchoice 16
\choice 24
\end{oneparchoices}

---
tipo: Múltipla Escolha
dificuldade: Média
nivel: médio
disciplina: Matemática
assunto: Análise Combinatória
gabarito: C
tags: [vôlei, duplas, canhotos, amigos]
fonte: concurso
concurso: ENEM
banca: INEP
ano: 2025
numero: 2
---
\question
Durante suas férias, oito amigos, dos quais dois são canhotos, decidem realizar um torneio de vôlei de praia. Eles precisam formar quatro duplas para a realização do torneio. Nenhuma dupla pode ser formada por dois jogadores canhotos.

De quantas maneiras diferentes podem ser formadas essas quatro duplas?
\begin{oneparchoices}
\choice 69
\choice 70
\correctchoice 90
\choice 104
\choice 105
\end{oneparchoices}

---
tipo: Múltipla Escolha
dificuldade: Fácil
nivel: médio
disciplina: Matemática
assunto: Análise Combinatória
gabarito: E
tags: [tênis, eliminatória, partidas, torneio]
fonte: concurso
concurso: ENEM
banca: INEP
ano: 2025
numero: 4
---
\question
Torneios de tênis, em geral, são disputados em sistema de eliminatória simples. Nesse sistema, são disputadas partidas entre dois competidores, com a eliminação do perdedor e promoção do vencedor para a fase seguinte. Dessa forma, se na 1ª fase o torneio conta com 2n competidores, então na 2ª fase restarão n competidores, e assim sucessivamente até a partida final.

Em um torneio de tênis, disputado nesse sistema, participam 128 tenistas.

Para se definir o campeão desse torneio, o número de partidas necessárias é dado por
\begin{oneparchoices}
\choice 2 × 128
\choice 64 + 32 + 16 + 8 + 4 + 2
\choice 128 + 64 + 32 + 16 + 8 + 4 + 2 + 1
\choice 128 + 64 + 32 + 16 + 8 + 4 + 2
\correctchoice 64 + 32 + 16 + 8 + 4 + 2 + 1
\end{oneparchoices}

---
tipo: Múltipla Escolha
dificuldade: Média
nivel: médio
disciplina: Matemática
assunto: Análise Combinatória
gabarito: E
tags: [piscina, produto, volume, tratamento]
fonte: concurso
concurso: ENEM
banca: INEP
ano: 2025
numero: 7
---
\question
Uma empresa especializada em conservação de piscinas utiliza um produto para tratamento da água cujas especificações técnicas sugerem que seja adicionado 1,5 mL desse produto para cada 1 000 L de água da piscina. Essa empresa foi contratada para cuidar de uma piscina de base retangular, de profundidade constante igual a 1,7 m, com largura e comprimento iguais a 3 m e 5 m, respectivamente. O nível da lâmina d’água dessa piscina é mantido a 50 cm da borda da piscina.

A quantidade desse produto, em mililitro, que deve ser adicionada a essa piscina de modo a atender às suas especificações técnicas é
\begin{oneparchoices}
\choice 11,25
\choice 27,00
\choice 28,80
\choice 32,25
\correctchoice 49,50
\end{oneparchoices}

---
tipo: Múltipla Escolha
dificuldade: Fácil
nivel: médio
disciplina: Matemática
assunto: Análise Combinatória
gabarito: C
tags: [símbolos, cores, daltonismo, Miguel Neiva]
fonte: concurso
concurso: ENEM
banca: INEP
ano: 2025
numero: 16
---
\question
O designer português Miguel Neiva criou um sistema de símbolos que permite que pessoas daltônicas identifiquem cores. O sistema consiste na utilização de símbolos que identificam as cores primárias (azul, amarelo e vermelho). Além disso, a justaposição de dois desses símbolos permite identificar cores secundárias, como o verde, que é o amarelo combinado com o azul. O preto e o branco são identificados por pequenos quadrados: o que simboliza o preto é cheio, enquanto o que simboliza o branco é vazio. Os símbolos que representam preto e branco também podem estar associados aos símbolos que identificam cores, significando se estas são claras ou escuras.

De acordo com o texto, quantas cores podem ser representadas pelo sistema proposto?
\begin{oneparchoices}
\choice 14
\choice 18
\correctchoice 20
\choice 21
\choice 23
\end{oneparchoices}

---
tipo: Múltipla Escolha
dificuldade: Fácil
nivel: médio
disciplina: Matemática
assunto: Análise Combinatória
gabarito: A
tags: [escola, brincadeira, respostas, objetos]
fonte: concurso
concurso: ENEM
banca: INEP
ano: 2025
numero: 17
---
\question
O diretor de uma escola convidou os 280 alunos de terceiro ano a participarem de uma brincadeira. Suponha que existem 5 objetos e 6 personagens numa casa de 9 cômodos; um dos personagens esconde um dos objetos em um dos cômodos da casa. O objetivo da brincadeira é adivinhar qual objeto foi escondido por qual personagem e em qual cômodo da casa o objeto foi escondido.

Todos os alunos decidiram participar. A cada vez um aluno é sorteado e dá a sua resposta. As respostas devem ser sempre distintas das anteriores, e um mesmo aluno não pode ser sorteado mais de uma vez. Se a resposta do aluno estiver correta, ele é declarado vencedor e a brincadeira é encerrada.

O diretor sabe que algum aluno acertará a resposta porque há
\begin{choices}
\correctchoice 10 alunos a mais do que possíveis respostas distintas.
\choice 20 alunos a mais do que possíveis respostas distintas.
\choice 119 alunos a mais do que possíveis respostas distintas.
\choice 260 alunos a mais do que possíveis respostas distintas.
\choice 270 alunos a mais do que possíveis respostas distintas.
\end{choices}

---
tipo: Múltipla Escolha
dificuldade: Média
nivel: médio
disciplina: Matemática
assunto: Análise Combinatória
gabarito: E
tags: [entrevista, números, algarismos, ordem]
fonte: concurso
concurso: ENEM
banca: INEP
ano: 2025
numero: 18
---
\question
O setor de recursos humanos de uma empresa vai realizar uma entrevista com 120 candidatos a uma vaga de contador. Por sorteio, eles pretendem atribuir a cada candidato um número, colocar a lista de números em ordem numérica crescente e usá-la para convocar os interessados. Acontece que, por um defeito do computador, foram gerados números com 5 algarismos distintos e, em nenhum deles, apareceram dígitos pares. Em razão disso, a ordem de chamada do candidato que tiver recebido o número 75 913 é
\begin{oneparchoices}
\choice 24
\choice 31
\choice 32
\choice 88
\correctchoice 89
\end{oneparchoices}

---
tipo: Múltipla Escolha
dificuldade: Fácil
nivel: médio
disciplina: Matemática
assunto: Análise Combinatória
gabarito: A
tags: [futebol, grupo A, abertura, mandante]
fonte: concurso
concurso: ENEM
banca: INEP
ano: 2025
numero: 21
---
\question
Doze times se inscreveram em um torneio de futebol amador. O jogo de abertura do torneio foi escolhido da seguinte forma: primeiro foram sorteados 4 times para compor o Grupo A. Em seguida, entre os times do Grupo A, foram sorteados 2 times para realizar o jogo de abertura do torneio, sendo que o primeiro deles jogaria em seu próprio campo, e o segundo seria o time visitante.

A quantidade total de escolhas possíveis para o Grupo A e a quantidade total de escolhas dos times do jogo de abertura podem ser calculadas através de
\begin{choices}
\correctchoice uma combinação e um arranjo, respectivamente.
\choice um arranjo e uma combinação, respectivamente.
\choice um arranjo e uma permutação, respectivamente.
\choice duas combinações.
\choice dois arranjos.
\end{choices}

---
tipo: Múltipla Escolha
dificuldade: Difícil
nivel: médio
disciplina: Matemática
assunto: Divisibilidade
gabarito: A
tags: [CPF, dígitos verificadores, cadastro, algoritmo]
fonte: concurso
concurso: ENEM
banca: INEP
ano: 2025
numero: 22
---
\question
Para cada indivíduo, a sua inscrição no Cadastro de Pessoas Físicas (CPF) é composta por um número de 9 algarismos e outro número de 2 algarismos, na forma \( d_1d_2 \), em que os dígitos \( d_1 \) e \( d_2 \) são denominados dígitos verificadores. Os dígitos verificadores são calculados, a partir da esquerda, da seguinte maneira: os 9 primeiros algarismos são multiplicados pela sequência 10, 9, 8, 7, 6, 5, 4, 3, 2; em seguida, calcula-se o resto \( r \) da divisão da soma dos resultados das multiplicações por 11, e se esse resto \( r \) for 0 ou 1, então \( d_1 = 0 \); caso contrário, \( d_1 = 11-r \). O dígito \( d_2 \) é obtido pela aplicação do mesmo processo, mas agora considerando os 9 primeiros algarismos seguidos de \( d_1 \), multiplicados pela sequência 11, 10, 9, 8, 7, 6, 5, 4, 3, 2.

Aplicando esse procedimento ao número 111.444.777, os dígitos verificadores obtidos são
\begin{oneparchoices}
\correctchoice 0 e 9
\choice 1 e 4
\choice 1 e 7
\choice 9 e 1
\choice 0 e 1
\end{oneparchoices}

---
tipo: Múltipla Escolha
dificuldade: Média
nivel: médio
disciplina: Matemática
assunto: Probabilidade Simples
gabarito: A
tags: [sorteio, colégio, inscrição, urnas]
fonte: concurso
concurso: ENEM
banca: INEP
ano: 2025
numero: 1
---
\question
Em um colégio público, a admissão no primeiro ano se dá por sorteio. Neste ano há 55 candidatos, cujas inscrições são numeradas de 01 a 55. O sorteio de cada número de inscrição será realizado em etapas, utilizando-se duas urnas. Da primeira urna será sorteada uma bola, dentre bolas numeradas de 0 a 9, que representará o algarismo das unidades do número de inscrição a ser sorteado e, em seguida, da segunda urna, será sorteada uma bola para representar o algarismo das dezenas desse número. Depois do primeiro sorteio, e antes de se sortear o algarismo das dezenas, as bolas que estarão presentes na segunda urna serão apenas aquelas cujos números formam, com o algarismo já sorteado, um número de 01 a 55.

As probabilidades de os candidatos de inscrição número 50 e 02 serem sorteados são, respectivamente,
\begin{choices}
\correctchoice \( \frac{1}{50} \) e \( \frac{1}{60} \)
\choice \( \frac{1}{50} \) e \( \frac{1}{50} \)
\choice \( \frac{1}{50} \) e \( \frac{1}{10} \)
\choice \( \frac{1}{55} \) e \( \frac{1}{54} \)
\choice \( \frac{1}{100} \) e \( \frac{1}{100} \)
\end{choices}

---
tipo: Múltipla Escolha
dificuldade: Média
nivel: médio
disciplina: Matemática
assunto: Probabilidade Simples
gabarito: C
tags: [promoção, voucher, urnas, bolinhas]
fonte: concurso
concurso: ENEM
banca: INEP
ano: 2025
numero: 2
---
\question
Visando atrair mais clientes, o gerente de uma loja anunciou uma promoção em que cada cliente que realizar uma compra pode ganhar um voucher para ser usado em sua próxima compra. Para ganhar seu voucher, o cliente precisa retirar, ao acaso, uma bolinha de dentro de cada uma das urnas A e B disponibilizadas pelo gerente, nas quais há apenas bolinhas pretas e brancas. Atualmente, a probabilidade de se escolher, ao acaso, uma bolinha preta na urna A é igual a 20% e a probabilidade de se escolher uma bolinha preta na urna B é 25%. Ganha o voucher o cliente que retirar duas bolinhas pretas, uma de cada urna.

Com o passar dos dias, o gerente percebeu que, para a promoção ser viável aos negócios, era preciso ajustar a probabilidade de acerto do cliente sem alterar a regra da promoção. Para isso, resolveu alterar a quantidade de bolinhas brancas na urna B de forma que a probabilidade de um cliente ganhar o voucher passasse a ser menor ou igual a 1%. Sabe-se que a urna B tem 4 bolinhas pretas e que, em ambas as urnas, todas as bolinhas têm a mesma probabilidade de serem retiradas.

Qual é o número mínimo de bolinhas brancas que o gerente deve adicionar à urna B?
\begin{oneparchoices}
\choice 20
\choice 60
\correctchoice 64
\choice 68
\choice 80
\end{oneparchoices}

---
tipo: Múltipla Escolha
dificuldade: Difícil
nivel: médio
disciplina: Matemática
assunto: Juros Compostos
gabarito: D
tags: [financiamento, parcelas, loja, taxa]
fonte: concurso
concurso: ENEM
banca: INEP
ano: 2025
numero: 3
---
\question
Uma loja vende seus produtos de duas formas: à vista ou financiado em três parcelas mensais iguais. Para definir o valor dessas parcelas nas vendas financiadas, a loja aumenta em 20% o valor do produto à vista e divide esse novo valor por 3. A primeira parcela deve ser paga no ato da compra, e as duas últimas, em 30 e 60 dias após a compra.

Um cliente da loja decidiu comprar, de forma financiada, um produto cujo valor à vista é R$ 1.500,00.

Utilize 5,29 como aproximação para \( \sqrt{28} \).

A taxa mensal de juros compostos praticada nesse financiamento é de
\begin{oneparchoices}
\choice 6,7%
\choice 10%
\choice 20%
\correctchoice 21,5%
\choice 23,3%
\end{oneparchoices}

---
tipo: Múltipla Escolha
dificuldade: Fácil
nivel: médio
disciplina: Matemática
assunto: Probabilidade Simples
gabarito: E
tags: [alojamento, estudantes, sorteio, cartões]
fonte: concurso
concurso: ENEM
banca: INEP
ano: 2025
numero: 4
---
\question
No alojamento de uma universidade, há alguns quartos com o padrão superior ao dos demais. Um desses quartos ficou disponível, e muitos estudantes se candidataram para morar no local. Para escolher quem ficará com o quarto, um sorteio será realizado. Para esse sorteio, cartões individuais com os nomes de todos os estudantes inscritos serão depositados em uma urna, sendo que, para cada estudante de primeiro ano, será depositado um único cartão com seu nome; para cada estudante de segundo ano, dois cartões com seu nome; e, para cada estudante de terceiro ano, três cartões com seu nome. Foram inscritos 200 estudantes de primeiro ano, 150 de segundo ano e 100 de terceiro ano. Todos os cartões têm a mesma probabilidade de serem sorteados.

Qual a probabilidade de o vencedor do sorteio ser um estudante de terceiro ano?
\begin{choices}
\choice \( \frac{1}{2} \)
\choice \( \frac{1}{3} \)
\choice \( \frac{1}{8} \)
\choice \( \frac{1}{9} \)
\correctchoice \( \frac{3}{8} \)
\end{choices}

---
tipo: Múltipla Escolha
dificuldade: Média
nivel: médio
disciplina: Matemática
assunto: Probabilidade Simples
gabarito: C
tags: [beisebol, World Series, campeão, partidas]
fonte: concurso
concurso: ENEM
banca: INEP
ano: 2025
numero: 5
---
\question
A World Series é a decisão do campeonato norte-americano de beisebol. Os dois times que chegam a essa fase jogam, entre si, até sete partidas. O primeiro desses times que completar quatro vitórias é declarado campeão.

Considere que, em todas as partidas, a probabilidade de qualquer um dos dois times vencer é sempre \( \frac{1}{2} \).

Qual é a probabilidade de o time campeão ser aquele que venceu a primeira partida da World Series?
\begin{choices}
\choice \( \frac{35}{64} \)
\choice \( \frac{40}{64} \)
\correctchoice \( \frac{42}{64} \)
\choice \( \frac{44}{64} \)
\choice \( \frac{52}{64} \)
\end{choices}

---
tipo: Múltipla Escolha
dificuldade: Média
nivel: médio
disciplina: Matemática
assunto: Probabilidade Simples
gabarito: D
tags: [dardos, campeonato, alvo, rodada]
fonte: concurso
concurso: ENEM
banca: INEP
ano: 2025
numero: 9
---
\question
O organizador de uma competição de lançamento de dardos pretende tornar o campeonato mais competitivo. Pelas regras atuais da competição, numa rodada, o jogador lança 3 dardos e pontua caso acerte pelo menos um deles no alvo. O organizador considera que, em média, os jogadores têm, em cada lançamento, \( \frac{1}{2} \) de probabilidade de acertar um dardo no alvo.

A fim de tornar o jogo mais atrativo, planeja modificar as regras de modo que a probabilidade de um jogador pontuar em uma rodada seja igual ou superior a \( \frac{9}{10} \). Para isso, decide aumentar a quantidade de dardos a serem lançados em cada rodada.

Com base nos valores considerados pelo organizador da competição, a quantidade mínima de dardos que devem ser disponibilizados em uma rodada para tornar o jogo mais atrativo é
\begin{oneparchoices}
\choice 2
\choice 4
\choice 6
\correctchoice 9
\choice 10
\end{oneparchoices}

---
tipo: Múltipla Escolha
dificuldade: Fácil
nivel: médio
disciplina: Matemática
assunto: Probabilidade Simples
gabarito: E
tags: [cofre, senha, dígitos, tentativa]
fonte: concurso
concurso: ENEM-PPL
banca: INEP
ano: 2021
numero: 10
---
\question
A senha de um cofre é uma sequência formada por oito dígitos, que são algarismos escolhidos de 0 a 9. Ao inseri-la, o usuário se esqueceu dos dois últimos dígitos que formam essa senha, lembrando somente que esses dígitos são distintos.

Digitando ao acaso os dois dígitos esquecidos, a probabilidade de que o usuário acerte a senha na primeira tentativa é
\begin{choices}
\choice \( \frac{2}{8} \)
\choice \( \frac{1}{90} \)
\choice \( \frac{2}{90} \)
\choice \( \frac{1}{100} \)
\correctchoice \( \frac{2}{100} \)
\end{choices}

---
tipo: Múltipla Escolha
dificuldade: Difícil
nivel: médio
disciplina: Matemática
assunto: Probabilidade Condicional
gabarito: C
tags: [fábrica, peças defeituosas, máquinas, produção]
fonte: concurso
concurso: ENEM-PPL
banca: INEP
ano: 2021
numero: 11
---
\question
Em uma fábrica de circuitos elétricos, há diversas linhas de produção e montagem. De acordo com o controle de qualidade da fábrica, as peças produzidas devem seguir um padrão. Em um processo produtivo, nem todas as peças produzidas são totalmente aproveitáveis, ou seja, há um percentual de peças defeituosas que são descartadas. Em uma linha de produção dessa fábrica, trabalham três máquinas, \( M_1 \), \( M_2 \) e \( M_3 \), dia e noite. A máquina \( M_1 \) produz 25% das peças, a máquina \( M_2 \) produz 30% e a máquina \( M_3 \) produz 45%. O percentual de peças defeituosas da máquina \( M_1 \) é de 2%, da máquina \( M_2 \) é de 3% e da máquina \( M_3 \) é igual a 4%.

A probabilidade de uma peça defeituosa ter sido produzida pela máquina \( M_2 \) é mais próxima de
\begin{oneparchoices}
\choice 15,6%
\choice 28,1%
\correctchoice 43,7%
\choice 56,2%
\choice 71,8%
\end{oneparchoices}

---
tipo: Múltipla Escolha
dificuldade: Difícil
nivel: médio
disciplina: Matemática
assunto: Probabilidade Condicional
gabarito: E
tags: [receita federal, fraude, inconsistência, declaração]
fonte: concurso
concurso: ENEM
banca: INEP
ano: 2025
numero: 22
---
\question
Em um determinado ano, os computadores da receita federal de um país identificaram como inconsistentes 20% das declarações de imposto de renda que lhe foram encaminhadas. Uma declaração é classificada como inconsistente quando apresenta algum tipo de erro ou conflito nas informações prestadas. Essas declarações consideradas inconsistentes foram analisadas pelos auditores, que constataram que 25% delas eram fraudulentas. Constatou-se ainda que, dentre as declarações que não apresentaram inconsistências, 6,25% eram fraudulentas.

Qual é a probabilidade de, nesse ano, a declaração de um contribuinte ser considerada inconsistente, dado que ela era fraudulenta?
\begin{oneparchoices}
\choice 0,0500
\choice 0,1000
\choice 0,1125
\choice 0,3125
\correctchoice 0,5000
\end{oneparchoices}

---
tipo: Múltipla Escolha
dificuldade: Difícil
nivel: médio
disciplina: Matemática
assunto: Probabilidade Simples
gabarito: B
tags: [prêmio, urnas, bolas pretas, escolha]
fonte: concurso
concurso: ENEM
banca: INEP
ano: 2025
numero: 23
---
\question
Para ganhar um prêmio, uma pessoa deverá retirar, sucessivamente e sem reposição, duas bolas pretas de uma mesma urna.

Inicialmente, as quantidades e cores das bolas são como descritas a seguir:

-- urna A: possui três bolas brancas, duas bolas pretas e uma bola verde;

-- urna B: possui seis bolas brancas, três bolas pretas e uma bola verde;

-- urna C: possui duas bolas pretas e duas bolas verdes;

-- urna D: possui três bolas brancas e três bolas pretas.

A pessoa deve escolher uma entre as cinco opções apresentadas:

-- opção 1: retirar, aleatoriamente, duas bolas da urna A;

-- opção 2: retirar, aleatoriamente, duas bolas da urna B;

-- opção 3: passar, aleatoriamente, uma bola da urna C para a urna A; após isso, retirar, aleatoriamente, duas bolas da urna A;

-- opção 4: passar, aleatoriamente, uma bola da urna D para a urna C; após isso, retirar, aleatoriamente, duas bolas da urna C;

-- opção 5: passar, aleatoriamente, uma bola da urna C para a urna D; após isso, retirar, aleatoriamente, duas bolas da urna D.

Com o objetivo de obter a maior probabilidade possível de ganhar o prêmio, a pessoa deve escolher a opção
\begin{oneparchoices}
\choice 1
\correctchoice 2
\choice 3
\choice 4
\choice 5
\end{oneparchoices}

---
tipo: Múltipla Escolha
dificuldade: Média
nivel: médio
disciplina: Matemática
assunto: Probabilidade Simples
gabarito: C
tags: [chuva, atraso, trabalho, meteorologia]
fonte: concurso
concurso: ENEM
banca: INEP
ano: 2017
numero: 1
---
\question
Um morador de uma região metropolitana tem 50% de probabilidade de atrasar-se para o trabalho quando chove na região; caso não chova, sua probabilidade de atraso é de 25%. Para um determinado dia, o serviço de meteorologia estima em 30% a probabilidade da ocorrência de chuva nessa região.

Qual é a probabilidade de esse morador se atrasar para o serviço no dia para o qual foi dada a estimativa de chuva?
\begin{oneparchoices}
\choice 0,075
\choice 0,150
\correctchoice 0,325
\choice 0,600
\choice 0,800
\end{oneparchoices}

---
tipo: Múltipla Escolha
dificuldade: Difícil
nivel: médio
disciplina: Matemática
assunto: Probabilidade Simples
gabarito: B
tags: [psicólogo, teste, verdadeiro ou falso, erro]
fonte: concurso
concurso: ENEM
banca: INEP
ano: 2014
numero: 4
---
\question
O psicólogo de uma empresa aplica um teste para analisar a aptidão de um candidato a determinado cargo. O teste consiste em uma série de perguntas cujas respostas devem ser verdadeiro ou falso e termina quando o psicólogo fizer a décima pergunta ou quando o candidato der a segunda resposta errada. Com base em testes anteriores, o psicólogo sabe que a probabilidade de o candidato errar uma resposta é 0,20.

A probabilidade de o teste terminar na quinta pergunta é
\begin{oneparchoices}
\choice 0,02048
\correctchoice 0,08192
\choice 0,24000
\choice 0,40960
\choice 0,49152
\end{oneparchoices}

---
tipo: Múltipla Escolha
dificuldade: Média
nivel: médio
disciplina: Matemática
assunto: Probabilidade Simples
gabarito: A
tags: [idiomas, escola, inglês, espanhol]
fonte: concurso
concurso: ENEM
banca: INEP
ano: 2013
numero: 7
---
\question
Numa escola com 1 200 alunos foi realizada uma pesquisa sobre o conhecimento desses em duas línguas estrangeiras, inglês e espanhol. Nessa pesquisa constatou-se que 600 alunos falam inglês, 500 falam espanhol e 300 não falam qualquer um desses idiomas.

Escolhendo-se um aluno dessa escola ao acaso e sabendo-se que ele não fala inglês, qual a probabilidade de que esse aluno fale espanhol?
\begin{choices}
\correctchoice \( \frac{1}{2} \)
\choice \( \frac{5}{8} \)
\choice \( \frac{1}{4} \)
\choice \( \frac{5}{6} \)
\choice \( \frac{5}{14} \)
\end{choices}

---
tipo: Múltipla Escolha
dificuldade: Fácil
nivel: médio
disciplina: Matemática
assunto: Probabilidade Simples
gabarito: D
tags: [dados, soma, jogo, jogadores]
fonte: concurso
concurso: ENEM
banca: INEP
ano: 2012
numero: 9
---
\question
José, Paulo e Antônio estão jogando dados não viciados, nos quais, em cada uma das seis faces, há um número de 1 a 6. Cada um deles jogará dois dados simultaneamente. José acredita que, após jogar seus dados, os números das faces voltadas para cima lhe darão uma soma igual a 7. Já Paulo acredita que sua soma será igual a 4 e Antônio acredita que sua soma será igual a 8.

Com essa escolha, quem tem a maior probabilidade de acertar sua respectiva soma é
\begin{choices}
\choice Antônio, já que sua soma é a maior de todas as escolhidas.
\choice José e Antônio, já que há 6 possibilidades tanto para a escolha de José quanto para a escolha de Antônio, e há apenas 4 possibilidades para a escolha de Paulo.
\choice José e Antônio, já que há 3 possibilidades tanto para a escolha de José quanto para a escolha de Antônio, e há apenas 2 possibilidades para a escolha de Paulo.
\correctchoice José, já que há 6 possibilidades para formar sua soma, 5 possibilidades para formar a soma de Antônio e apenas 3 possibilidades para formar a soma de Paulo.
\choice Paulo, já que sua soma é a menor de todas.
\end{choices}

---
tipo: Múltipla Escolha
dificuldade: Média
nivel: médio
disciplina: Matemática
assunto: Probabilidade Simples
gabarito: C
tags: [sinuca, bolas, soma, caçapa]
fonte: concurso
concurso: ENEM
banca: INEP
ano: 2011
numero: 13
---
\question
Em um jogo disputado em uma mesa de sinuca, há 16 bolas: 1 branca e 15 coloridas, as quais, de acordo com a coloração, valem de 1 a 15 pontos, um valor para cada bola colorida. O jogador acerta o taco na bola branca de forma que esta acerte as outras, com o objetivo de acertar duas das quinze bolas em quaisquer caçapas. Os valores dessas duas bolas são somados e devem resultar em um valor escolhido pelo jogador antes do início da jogada.

Arthur, Bernardo e Caio escolhem os números 12, 17 e 22 como sendo resultados de suas respectivas somas.

Com essa escolha, quem tem a maior probabilidade de ganhar é
\begin{choices}
\choice Arthur, pois a soma que escolheu é a menor.
\choice Bernardo, pois há 7 possibilidades de compor a soma escolhida por ele, contra 4 possibilidades para a escolha de Arthur e 4 possibilidades para a escolha de Caio.
\correctchoice Bernardo, pois há 7 possibilidades de compor a soma escolhida por ele, contra 5 possibilidades para a escolha de Arthur e 4 possibilidades para a escolha de Caio.
\choice Caio, pois há 10 possibilidades de compor a soma escolhida por ele, contra 5 possibilidades para a escolha de Arthur e 8 possibilidades para a escolha de Bernardo.
\choice Caio, pois a soma que escolheu é a maior.
\end{choices}

---
tipo: Múltipla Escolha
dificuldade: Média
nivel: médio
disciplina: Matemática
assunto: Probabilidade Simples
gabarito: C
tags: [celular, defeito, loja, aparelhos]
fonte: concurso
concurso: ENEM
banca: INEP
ano: 2009
numero: 18
---
\question
O controle de qualidade de uma empresa fabricante de telefones celulares aponta que a probabilidade de um aparelho de determinado modelo apresentar defeito de fabricação é de 0,2%. Se uma loja acaba de vender 4 aparelhos desse modelo para um cliente, qual é a probabilidade de esse cliente sair da loja com exatamente dois aparelhos defeituosos?
\begin{choices}
\choice \( 2 \cdot (0,2\%)^4 \)
\choice \( 4 \cdot (0,2\%)^2 \)
\correctchoice \( 6 \cdot (0,2\%)^2 \cdot (99,8\%)^2 \)
\choice \( 4 \cdot (0,2\%) \)
\choice \( 6 \cdot (0,2\%) \cdot (99,8\%) \)
\end{choices}

---
tipo: Múltipla Escolha
dificuldade: Difícil
nivel: médio
disciplina: Matemática
assunto: Probabilidade Simples
gabarito: C
tags: [mega-sena, quina, apostas, loteria]
fonte: concurso
concurso: ENEM
banca: INEP
ano: 2009
numero: 20
---
\question
A população brasileira sabe, pelo menos intuitivamente, que a probabilidade de acertar as seis dezenas da mega sena não é zero, mas é quase. Mesmo assim, milhões de pessoas são atraídas por essa loteria, especialmente quando o prêmio se acumula em valores altos. Até junho de 2009, cada aposta de seis dezenas, pertencentes ao conjunto \{01, 02, 03, ..., 59, 60\}, custava R$ 1,50.

Considere que uma pessoa decida apostar exatamente R$ 126,00 e que esteja mais interessada em acertar apenas cinco das seis dezenas da mega sena, justamente pela dificuldade desta última. Nesse caso, é melhor que essa pessoa faça 84 apostas de seis dezenas diferentes, que não tenham cinco números em comum, do que uma única aposta com nove dezenas, porque a probabilidade de acertar a quina no segundo caso em relação ao primeiro é, aproximadamente,
\begin{choices}
\choice \( 1\frac{1}{2} \)
\choice \( 2\frac{1}{2} \)
\correctchoice 4 vezes menor.
\choice 9 vezes menor.
\choice 14 vezes menor.
\end{choices}

---
tipo: Múltipla Escolha
dificuldade: Média
nivel: médio
disciplina: Matemática
assunto: Combinações
gabarito: A
tags: [tênis, exibição, canhotos, destros]
fonte: concurso
concurso: ENEM
banca: INEP
ano: 2025
numero: 12
---
\question
O tênis é um esporte em que a estratégia de jogo a ser adotada depende, entre outros fatores, de o adversário ser canhoto ou destro.

Um clube tem um grupo de 10 tenistas, sendo que 4 são canhotos e 6 são destros. O técnico do clube deseja realizar uma partida de exibição entre dois desses jogadores, porém, não poderão ser ambos canhotos.

Qual o número de possibilidades de escolha dos tenistas para a partida de exibição?
\begin{choices}
\CorrectChoice \( \frac{10!}{2! \cdot 8!} - \frac{4!}{2! \cdot 2!} \)
\choice \( \frac{10!}{8!} - \frac{4!}{2!} \)
\choice \( \frac{10!}{2! \cdot 8!} - 2 \)
\choice \( \frac{6!}{4!} + 4 \times 4 \)
\choice \( \frac{6!}{4!} + 6 \times 4 \)
\end{choices}

---
tipo: Múltipla Escolha
dificuldade: Média
nivel: médio
disciplina: Matemática
assunto: Análise Combinatória
gabarito: E
tags: [senha, cadastro, site, caracteres]
fonte: concurso
concurso: ENEM
banca: INEP
ano: 2025
numero: 13
---
\question
Para cadastrar-se em um site, uma pessoa precisa escolher uma senha composta por quatro caracteres, sendo dois algarismos e duas letras (maiúsculas ou minúsculas). As letras e os algarismos podem estar em qualquer posição. Essa pessoa sabe que o alfabeto é composto por vinte e seis letras e que uma letra maiúscula difere da minúscula em uma senha.

Disponível em: infowester.com. Acesso em: 14 dez. 2012.

O número total de senhas possíveis para o cadastramento nesse site é dado por
\begin{choices}
\choice \( 10^2 \cdot 26^2 \)
\choice \( 10^2 \cdot 52^2 \)
\choice \( 10^2 \cdot 52^2 \cdot \frac{4!}{2!} \)
\choice \( 10^2 \cdot 26^2 \cdot \frac{4!}{2! \cdot 2!} \)
\CorrectChoice \( 10^2 \cdot 52^2 \cdot \frac{4!}{2! \cdot 2!} \)
\end{choices}

---
tipo: Múltipla Escolha
dificuldade: Média
nivel: médio
disciplina: Matemática
assunto: Permutações Simples
gabarito: B
tags: [videolocadora, filmes, ação, comédia]
fonte: concurso
concurso: ENEM
banca: INEP
ano: 2025
numero: 14
---
\question
Um cliente de uma videolocadora tem o hábito de alugar dois filmes por vez. Quando os devolve, sempre pega outros dois filmes e assim sucessivamente. Ele soube que a videolocadora recebeu alguns lançamentos, sendo 8 filmes de ação, 5 de comédia e 3 de drama e, por isso, estabeleceu uma estratégia para ver todos esses 16 lançamentos. Inicialmente alugará, em cada vez, um filme de ação e um de comédia. Quando se esgotarem as possibilidades de comédia, o cliente alugará um filme de ação e um de drama, até que todos os lançamentos sejam vistos e sem que nenhum filme seja repetido.

De quantas formas distintas a estratégia desse cliente poderá ser posta em prática?
\begin{choices}
\choice \( 20 \times 8! + (3!)^2 \)
\CorrectChoice \( 8! \times 5! \times 3! \)
\choice \( \frac{8! \times 5! \times 3!}{2^8} \)
\choice \( \frac{8! \times 5! \times 3!}{2^2} \)
\choice \( \frac{16!}{2^8} \)
\end{choices}

---
tipo: Múltipla Escolha
dificuldade: Média
nivel: médio
disciplina: Matemática
assunto: Estatística
gabarito: D
tags: [ônibus, mediana, moda, horário]
fonte: concurso
concurso: ENEM
banca: INEP
ano: 2025
numero: 25
---
\question
Um rapaz estuda em uma escola que fica longe de sua casa, e por isso precisa utilizar o transporte público. Como é muito observador, todos os dias ele anota a hora exata (sem considerar os segundos) em que o ônibus passa pelo ponto de espera.

Também notou que nunca consegue chegar ao ponto de ônibus antes de 6h 15min da manhã. Analisando os dados coletados durante o mês de fevereiro, o qual teve 21 dias letivos, ele concluiu que 6h 21min foi o que mais se repetiu, e que a mediana do conjunto de dados é 6h 22min.

A probabilidade de que, em algum dos dias letivos de fevereiro, esse rapaz tenha apanhado o ônibus antes de 6h 21min da manhã é, no máximo,
\begin{choices}
\choice \( \frac{4}{21} \)
\choice \( \frac{5}{21} \)
\choice \( \frac{6}{21} \)
\CorrectChoice \( \frac{7}{21} \)
\choice \( \frac{8}{21} \)
\end{choices}
